\documentclass[10pt,a4paper]{article}
\usepackage[utf8]{inputenc}
\usepackage[T1]{fontenc}
\usepackage[brazilian]{babel}
\usepackage{graphicx}
\usepackage{amsmath}
\usepackage{booktabs}
\usepackage[margin=2.2cm]{geometry}
\usepackage{hyperref}
\usepackage{float}

% Mostra o gráfico se existir; senão, uma caixa de aviso (não quebra a compilação).
\newcommand{\grafico}[2]{%
  \begin{figure}[H]\centering
  \IfFileExists{../resultados/graficos/#1.png}%
    {\includegraphics[width=0.82\linewidth]{../resultados/graficos/#1.png}}%
    {\fbox{\parbox{0.8\linewidth}{\centering [gráfico \texttt{\detokenize{#1}} ainda não gerado]}}}%
  \caption{#2}\label{fig:#1}
  \end{figure}}
\newcommand{\preencher}[1]{\textbf{[PREENCHER: #1]}}

\title{Otimização do Método de Newton para Sistemas Não-Lineares de Broyden:\\ análise de desempenho}
\author{Nome1 (GRR) \and Nome2 (GRR)\\
Departamento de Informática -- UFPR\\ CI1164 -- Introdução à Computação Científica}
\date{\today}

\begin{document}
\maketitle

\begin{abstract}
Este trabalho implementa o Método de Newton para resolver Sistemas Não-Lineares de Broyden (SNLB)
e compara quatro versões do programa: uma base com matriz densa e solver $O(n^3)$, uma versão que explora a
estrutura tridiagonal da Jacobiana e duas versões otimizadas (vetorização, alinhamento, fusão de laços e
redução de divisões). O desempenho é avaliado com o LIKWID (tempo, banda de memória, taxa de cache miss
e MFLOP/s) para $N$ de 32 a 20000. \preencher{resumir os principais resultados numéricos (ex.: speedup da v3 sobre a v0 e a v1).}
\end{abstract}

\section{Introdução}
Um SNLB é $F(X)=0$, com $f_1=-2x_1^2+3x_1-2x_2+1$, $f_i=-2x_i^2+3x_i-x_{i-1}-2x_{i+1}+1$ ($2\le i\le n-1$)
e $f_n=-2x_n^2+3x_n-x_{n-1}$. O Método de Newton resolve, a cada iteração $k$, o sistema linear
$J(X^{(k)})\Delta^{(k)}=-F(X^{(k)})$ e atualiza $X^{(k+1)}=X^{(k)}+\Delta^{(k)}$, parando quando
$\|F(X^{(k)})\|_\infty<\varepsilon$, quando $\|\Delta^{(k)}\|_\infty<\varepsilon$ ou após MAX iterações.
Neste problema a Jacobiana é \emph{tridiagonal}: $\partial f_i/\partial x_{i-1}=-1$,
$\partial f_i/\partial x_i=3-4x_i$ e $\partial f_i/\partial x_{i+1}=-2$.

\section{Ambiente experimental}
\preencher{colar o resultado de \texttt{likwid-topology -g -c} (modelo da CPU, nº de núcleos, frequência, tamanhos e associatividade das caches L1/L2/L3, banda de memória). O arquivo gerado é \texttt{resultados/topologia.txt}.}
\par\smallskip
Compilador: GCC \preencher{versão}, flags \texttt{-O3 -march=native -mavx -fopt-info-vec}. Sistema: \preencher{SO/kernel}.
Frequência fixada em \preencher{GHz}, execução no núcleo \preencher{n} com a máquina exclusiva.
Todos os testes usam $\varepsilon=0$, MAX=25 e $x_0=0$; a versão densa é limitada a $N\leq 2000$ pelo custo $O(n^3)$, enquanto as versões tridiagonais são testadas até $N=20000$. Os tempos são a média de \preencher{10} execuções
(\texttt{timestamp()} em ms, binário sem instrumentação; a região NEWTON cobre o laço iterativo, sem alocações/inicialização nem impressão) e os contadores vêm de uma execução sob \texttt{likwid-perfctr}
com marcadores \texttt{NEWTON}, \texttt{JACOBIANA} e \texttt{SISTEMA\_LINEAR}.

\section{Implementação}
\subsection{Decisões de projeto}
\paragraph{Método do sistema linear.} A Jacobiana é tridiagonal, então a eliminação de Gauss
custa $O(n)$ (cerca de $5n$ FLOPs na triangularização e $3n$ na retrossubstituição, slides da Parte 5)
contra $O(n^3)$ de uma eliminação densa. Mantém-se pivoteamento parcial entre linhas adjacentes, pois $3-4x_i$
pode se aproximar de zero; o pivoteamento pode gerar uma segunda superdiagonal (\texttt{du2}).
\paragraph{Estrutura de dados.} Quatro vetores contíguos (\texttt{a,d,c,du2}) em vez de uma matriz $n\times n$:
memória $O(n)$ e acessos de passo unitário. Na versão base usa-se \texttt{double **} com linhas separadas.
\paragraph{Cálculo da Jacobiana.} Derivadas analíticas escritas diretamente (sem derivação numérica nem
\texttt{libmatheval}), custo $O(n)$ por iteração.

\subsection{Versões}
\begin{itemize}
\item \textbf{v0 (base)}: Jacobiana densa, Gauss $O(n^3)$, sem otimizações.
\item \textbf{v1}: tridiagonal, pivoteamento adjacente; \texttt{calloc/free} por iteração, três passadas de \texttt{isfinite}, vetores auxiliares.
\item \textbf{v2}: v1 mais as otimizações abaixo.
\item \textbf{v3}: v2 com um único inverso de pivô por linha.
\end{itemize}

\subsection{Otimizações e justificativas}
\begin{description}
\item[\texttt{restrict} e ausência de aliasing] (\texttt{funcao.c}, \texttt{jacobiana.c}, \texttt{solver.c}, \texttt{newton.c}).
Sem \texttt{restrict} o GCC versiona os laços por possível aliasing; com ele vetoriza direto (relatório \texttt{-fopt-info-vec}).
\item[Laços vetorizáveis] \texttt{calculaMenosF} e \texttt{calculaJacobiana} são laços de passo unitário, sem desvios e sem
dependência de dados, vetorizados com SIMD; a quantidade de \texttt{double} processada por instrução depende da ISA disponível na CPU (por exemplo, 4 em vetores de 256 bits e 8 em vetores de 512 bits).
\item[Eliminação de subexpressões e menos trabalho] $-f_i=x_i(2x_i-3)+x_{i-1}+2x_{i+1}-1$; calcula-se diretamente $b=-F$, eliminando um vetor e uma passada.
\item[Subdiagonal constante] \texttt{a[]} vale $-1$ e não é alterada pelo solver: é preenchida uma vez; a Jacobiana só reescreve \texttt{d} e \texttt{c}.
\item[Fusão de laços] norma e checagem de finitude de $b$ em uma passada; $X+\Delta$, $\|\Delta\|_\infty$ e finitude em outra (antes: 3 passadas). Reduz tráfego de memória, pois o código é limitado por memória.
\item[Sem alocação no laço] \texttt{du2} pré-alocado; elimina \texttt{calloc/free} (e a página zerada) no trecho \texttt{SISTEMA\_LINEAR}.
\item[Menos desvios] sem \texttt{isfinite} dentro da fatoração, uma só passada para a escala, sem testes de pivô repetidos na retrossubstituição.
\item[Alinhamento] \texttt{aligned\_alloc(64)} nos vetores de trabalho, favorecendo alinhamento com linhas de cache e instruções vetoriais; a ISA efetivamente usada depende da CPU.
\item[Inverso do pivô (v3)] a divisão tem latência alta e pouco pipeline (slides de pipelining); como a retrossubstituição é uma cadeia de dependências, substituir $/d_i$ por $\cdot(1/d_i)$ remove a divisão do caminho crítico.
\end{description}
\paragraph{O que não foi aplicado, e por quê.} Padding contra potência de 2 e \emph{loop blocking}/\emph{unroll\&jam}
não se aplicam: todos os acessos têm passo unitário e os dados são $O(n)$. A eliminação é intrinsecamente sequencial
(dependência de dados entre linhas), por isso não vetoriza.

\section{Resultados}
\subsection{Tempo de execução}
\grafico{tempo_newton}{Tempo médio -- execução total do Método de Newton.}
\grafico{tempo_jacobiana}{Tempo médio -- cálculo das matrizes Jacobianas.}
\grafico{tempo_sistema_linear}{Tempo médio -- resolução dos sistemas lineares.}
\preencher{comparar as curvas com as complexidades teóricas ($O(n^3)$ na v0 e $O(n)$ nas versões tridiagonais), calcular os speedups observados, discutir ruído em N pequeno e possíveis efeitos das caches em N grande.}

\subsection{Banda de memória}
\grafico{banda_newton}{Banda de memória -- execução total.}
\grafico{banda_jacobiana}{Banda de memória -- Jacobianas.}
\grafico{banda_sistema_linear}{Banda de memória -- sistemas lineares.}
\preencher{explicar o que mostram (grupo MEM ou L3), relacionar com o tamanho dos dados e das caches da CPU.}

\subsection{Cache miss}
\grafico{cache_newton}{Taxa de cache miss -- execução total.}
\grafico{cache_jacobiana}{Taxa de cache miss -- Jacobianas.}
\grafico{cache_sistema_linear}{Taxa de cache miss -- sistemas lineares.}
\preencher{identificar o grupo efetivamente usado (L2CACHE ou fallback), comparar as taxas observadas e relacioná-las aos layouts de memória e padrões de acesso; não presumir qual versão apresenta mais misses.}

\subsection{Operações aritméticas}
\grafico{flops_newton}{MFLOP/s -- execução total (FLOPS\_DP contínuo, FLOPS\_AVX tracejado).}
\grafico{flops_jacobiana}{MFLOP/s -- Jacobianas.}
\grafico{flops_sistema_linear}{MFLOP/s -- sistemas lineares.}
\preencher{analisar os valores de FLOPS\_DP e FLOPS\_AVX em cada região, relacionar com a vetorização observada e discutir possíveis limites de memória ou dependências sequenciais quando sustentados pelos dados.}

\section{Conclusão}
\preencher{conclusões, limitações (ex.: dependência de $x_0$ para a convergência) e trabalhos futuros.}

\begin{thebibliography}{9}
\bibitem{hager} G. Hager, G. Wellein. \emph{Introduction to High Performance Computing for Scientists and Engineers}. CRC Press, 2011.
\bibitem{ruggiero} M. A. G. Ruggiero, V. L. R. Lopes. \emph{Cálculo Numérico: Aspectos Teóricos e Computacionais}. McGraw-Hill, 1988.
\bibitem{peters} S. Peters, J. F. Szeremeta. \emph{Cálculo Numérico Computacional}. Editora UFSC.
\bibitem{likwid} \emph{LIKWID}. \url{https://github.com/RRZE-HPC/likwid}
\bibitem{slides} A. Delgado, G. Derenievicz. Notas de aula de CI1164 -- Introdução à Computação Científica, UFPR/DINF.
\end{thebibliography}
\end{document}
